\documentclass[11pt]{article}
\usepackage[margin=2.6cm]{geometry}
\usepackage{mathpazo}
\usepackage{microtype}
\usepackage{booktabs}
\usepackage{graphicx}
\usepackage{float}
\usepackage{xcolor}
\usepackage{titlesec}
\usepackage{caption}
\usepackage{listings}
\usepackage[hidelinks]{hyperref}
\emergencystretch=3em

\definecolor{accent}{HTML}{1F4E79}
\definecolor{good}{HTML}{1A7F37}
\definecolor{bad}{HTML}{B42318}
\titleformat{\section}{\Large\bfseries\color{accent}}{\thesection}{0.8em}{}
\titleformat{\subsection}{\large\bfseries\color{accent!80!black}}{\thesubsection}{0.8em}{}
\captionsetup{font=small,labelfont={bf,color=accent}}
\newcommand{\armA}{\textcolor{bad}{arm A}}
\newcommand{\armB}{\textcolor{accent}{arm B}}

\lstdefinestyle{json}{
  basicstyle=\ttfamily\footnotesize,
  backgroundcolor=\color{gray!8},
  frame=single, rulecolor=\color{gray!40}, framerule=0.4pt,
  breaklines=true, columns=fullflexible, keepspaces=true,
  showstringspaces=false, aboveskip=8pt, belowskip=8pt}
\lstset{style=json}

\newsavebox{\eliBox}
\newenvironment{eli5}
  {\par\medskip\noindent\begin{lrbox}{\eliBox}\begin{minipage}{\dimexpr\linewidth-2\fboxsep-2\fboxrule}\small}
  {\end{minipage}\end{lrbox}\noindent\fcolorbox{accent!40}{accent!4}{\usebox{\eliBox}}\par\medskip}

\title{\color{accent}\bfseries Tool-Aware DAG-Diff Mutation for Retrieval Chains\\
\large Strong-Mutator A/B on HoVer multi-hop retrieval: does the structured diff hold\\
when the genome carries \texttt{tool} steps and the mutator is a capable reasoning model?\\
\large (single-parent, \texttt{gemini-3.5-flash}, full 250-mutant runs)}
\author{GigaEvo}
\date{2026-07-04}

\begin{document}
\maketitle

\begin{abstract}
\noindent
This is the third A/B in the structured-DAG-diff line and the first on a task whose genome
carries \texttt{tool}/\texttt{retrieve} steps. We battle-test a \emph{tool-aware} slot-diff
grammar (\texttt{AllowedToolChainChanges}) against stock free-form wire-JSON rewriting on
\textbf{HoVer multi-hop retrieval} (\texttt{problems/chains/hover/full7}, full-chain mode:
variable topology, LLM + \texttt{retrieve}/\texttt{retrieve\_deep} tool steps), scoring each
child by soft retrieval coverage (fraction of gold supporting-fact titles found, 300 samples,
Qwen3-8B executor). Unlike the prior weak-\texttt{llama} run, the mutator here is a
\emph{strong reasoning model} (\texttt{google/gemini-3.5-flash}, \texttt{reasoning.effort:high});
both arms ran the full 250 mutants from the same 7-step seed with \texttt{num\_parents=1},
\texttt{memory=none}. The validity result is decisive and, notably, holds \emph{even for a
strong model}: over completed evaluations \armA{} (free rewrite) yielded \textbf{149/242 valid
(61.6\%)} with \textbf{91 malformed-JSON children (38\%)}, while \armB{} (tool diff) yielded
\textbf{239/242 valid (98.8\%)} with \textbf{zero malformed JSON} (3 benign non-parse
rejections). That gemini's malformed rate (38\%) is comparable to the weak model's (40\%) says
the failure is \emph{structural to full-document re-emission of a tool-bearing chain}, not
model weakness. \armB{} also wins on fitness --- child mean \textbf{0.800 vs 0.778}, best
\textbf{0.858 vs 0.831} --- and keeps chains tighter to the seed topology (84\% at 7 steps vs
77\%). The token story is the honest twist: on this run the diff is \emph{not} cheaper in total
output. Both arms requested identical reasoning settings, but the run's token log recorded a
reasoning channel only for the diff calls (87\% of \armB's output); \armA's reasoning was
\emph{never recorded} --- the \texttt{MutationAgent} \texttt{LLM\_CALL} emitter dropped the
\texttt{tokens\_reasoning} field, so \armA{} is an \emph{unmeasured}-reasoning arm, not a
0-reasoning one (bug since fixed; this run's logs are unrecoverable). The clean read that
survives is \emph{total} billed output (\texttt{tokens\_out}, correct on both arms): \armB's is
\textbf{+12\%} higher. Finally, the
experiment's engineering deliverable landed: per-candidate CARL executor token accounting
(\texttt{problems/chains/usage.py}) is now wired into \texttt{full7/validate.py}, closing the
gap that forced this run's executor cost to be reported only as a shared-proxy aggregate.
\end{abstract}

\section{Motivation}
CARL chain genomes are wire-JSON documents (\texttt{ReasoningChain.to\_dict} plus platform
extras). The stock mutation path (\texttt{mutation=llm\_rewrite}) asks the LLM to re-emit the
\emph{whole} document as free text: every emission re-risks the entire syntax, and a malformed
child is only discovered \emph{after} it has been persisted and (if the parse survives
extraction) has consumed a full evaluation. Two prior A/Bs established the structured-diff
alternative on the \emph{summarizer} task, whose chains are LLM-only. This run answers the two
questions those left open:
\begin{itemize}
  \item \textbf{Does the diff generalise to tool-bearing genomes?} HoVer chains interleave
  \texttt{retrieve}/\texttt{retrieve\_deep} \texttt{tool} steps with LLM steps. The original
  \texttt{ChainDagDiff} was LLM-step-only and hard-rejected tool steps; \armB{} exercises the
  new \texttt{AllowedToolChainChanges} vocabulary that admits tool steps as first-class slots.
  \item \textbf{Does the guarantee still matter when the mutator is \emph{strong}?} The prior
  run used a weak non-reasoning 8B model, where free-form JSON was expected to break. Here the
  mutator is \texttt{gemini-3.5-flash} with \texttt{reasoning.effort:high} --- a capable model
  that should rarely emit malformed JSON. If free rewriting still fails often here, the failure
  is structural, not a weak-model artifact.
\end{itemize}

\section{Setup}
\label{sec:setup}
Two arms, run from scratch, identical except the mutation operator (\texttt{launch\_arm.sh}):
\begin{itemize}
  \item \textbf{Task:} \texttt{problems/chains/hover/full7}, \texttt{pipeline=program\_is\_json}.
  Full-chain mode: no frozen steps, no topology matching; the child may freely change step
  count/types/dependencies within \texttt{FULL\_CHAIN\_CONFIG} (\texttt{max\_steps}=7). Fitness
  is soft retrieval coverage --- for each of 300 HoVer claims, the fraction of gold
  supporting-fact article titles found across \emph{all} tool-step outputs (higher better).
  \item \textbf{Mutator:} \texttt{google/gemini-3.5-flash} via OpenRouter
  (\texttt{llm=gemini35\_flash}, \texttt{reasoning.effort:high}, \texttt{service\_tier:flex})
  --- \emph{identical config on both arms}.
  \item \textbf{Executor:} \texttt{Qwen/Qwen3-8B} via the LiteLLM proxy at
  \texttt{10.232.89.98:4000}, run step-batched over the 300-sample set.
  \item \textbf{Config:} \texttt{memory=none}, \texttt{num\_parents=1}, one 7-step retrieval
  seed, disk storage, \texttt{max\_mutants=250}. \armA{} = \texttt{mutation=llm\_rewrite}
  (stock full wire-JSON rewrite); \armB{} = \texttt{mutation=carl\_with\_retrieval\_tools}
  (\texttt{StructuredDiffMutationOperator} + \texttt{AllowedToolChainChanges}).
  \item \textbf{Matching:} both arms ran the \emph{full} 250 mutants (no early stop, no
  iteration cap). Each arm had exactly \textbf{14 children still executing} at the 250-mutant
  stop; those in-flight children are excluded from the ``completed'' validity/fitness figures
  and noted where relevant.
\end{itemize}

\noindent\textbf{Single parent.} With \texttt{num\_parents=1} the diff grammar collapses to its
single-parent form (\texttt{base\_parent} is one letter; every \texttt{keep}-id enum ranges
over one parent's steps). Crossover is unavailable this run --- a deliberate simplification,
two-parent tool-crossover left as prior work.

\section{The tool-aware diff language}
\label{sec:language}
The child chain is a shelf of numbered slots; the LLM fills them left-to-right. What is new
here versus the summarizer diff is that a slot may hold a \texttt{tool} step, not only an LLM
step: a kept tool step carries its tool name and arguments, and the grammar admits editing them.

\begin{eli5}
Picture the child chain as a row of numbered boxes. Into each box you put either a
\emph{photocopy} of a step from the parent --- named by its label (\texttt{a1}, \texttt{a2},
\dots), optionally with small corrections --- or a \emph{brand-new} step. Some parent steps are
``go-look-it-up'' \emph{tool} steps (retrieve passages); those are photocopied and corrected the
same way as thinking steps. Each box lists which \emph{earlier} boxes it reads from, and the
form for box $k$ only has checkboxes for boxes $1..k{-}1$ --- so pointing forward or at yourself
is not against the rules, there is simply nowhere on the form to write it. Leftover boxes stay
empty (no gaps), and the last filled box produces the graded answer.
\end{eli5}
\noindent
The form is a JSON schema compiled fresh per mutation from the actual parent and shipped to the
provider as the \texttt{json\_schema} structured-output grammar, so a violating token can never
be sampled; a pydantic model re-checks as a backstop, and the applier round-trips the result
through \texttt{ReasoningChain} as a final assertion. Everything an unconstrained rewriter
routinely gets wrong --- self/forward dependencies, cycles, keeping a nonexistent step,
over-long chains, invented fields, \emph{malformed JSON} --- is unrepresentable by construction;
only step-contiguity is left to a validator. A rejection costs one LLM call: nothing persisted,
nothing evaluated.

\section{Results}
\label{sec:results}

\subsection{Validity funnel --- the headline}
\begin{table}[H]\centering
\begin{tabular}{lcc}
\toprule
full 250-mutant run & \armA{} (rewrite) & \armB{} (tool diff) \\
\midrule
mutation LLM calls (incl.\ retries) & 252 & 259 \\
pre-persist generation failures & 1 (no code) & 4 (2 schema, 2 call) \\
children persisted & 250 & 254 \\
\quad still executing at stop (excluded) & 14 & 14 \\
\midrule
\textbf{completed evaluations} & \textbf{242} & \textbf{242} \\
\quad valid & 149 & \textbf{239} \\
\quad invalid & 93 & \textbf{3} \\
\qquad of which malformed JSON & \textbf{\textcolor{bad}{91 (38\%)}} & \textbf{\textcolor{good}{0}} \\
\qquad other (non-parse) & 2 & 3 \\
\textbf{valid yield among completed} & \textbf{\textcolor{bad}{61.6\%}} & \textbf{\textcolor{good}{98.8\%}} \\
\bottomrule
\end{tabular}
\caption{The free-form path breaks on \emph{completed} children even with a strong mutator:
\textbf{91 of 242 (38\%) are malformed wire-JSON} (\texttt{json\_parse\_error}, decoded from the
parse stage), each having burned an evaluation before being marked invalid. \armB's tool diff
persisted \textbf{zero} malformed-JSON children; its 3 completed-invalid children failed a
downstream non-parse check, and its pre-persist rejections (2 diff-schema, 2 LLM-call) fired
before anything was persisted. Both arms left 14 children mid-evaluation at the 250-mutant stop;
those are excluded here (they are the bulk of the raw ``invalid'' count in
\texttt{ab\_stats.json}: \armB's 15 raw invalids are 12 running + 3 genuine).}
\end{table}
\begin{figure}[H]\centering
\includegraphics[width=\textwidth]{ab_overview.png}
\caption{Left: per-child soft-coverage fitness (dots) and best-so-far (line); invalid children
are drawn as crosses on the floor line (dense red for \armA, nearly absent for \armB). \armB's
valid cloud sits visibly higher and its best-so-far climbs to 0.858 vs \armA's 0.831. Right: the
validity funnel over the full run --- both persist $\sim$250, but \armA{} keeps 149 valid while
\armB{} keeps 239.}
\end{figure}

\noindent\textbf{Why 38\% for a strong model.} The weak-\texttt{llama} summarizer run malformed
$\sim$40\% of its rewrites; gemini-3.5-flash malforms 38\% here. A capable reasoning model is
\emph{not} materially better at re-emitting a whole tool-bearing chain as valid JSON --- the
defect scales with document complexity (nested tool arguments, dependency lists, platform
extras), not with model IQ. This is the cleanest evidence yet that the guarantee is worth its
schema cost: the diff removes a failure class that strong models do not outgrow.

\subsection{Fitness --- \armB{} wins mean and best}
\begin{table}[H]\centering
\begin{tabular}{lcc}
\toprule
completed children & \armA{} (rewrite) & \armB{} (tool diff) \\
\midrule
best child & 0.831 & \textbf{\textcolor{good}{0.858}} \\
child mean & 0.778 & \textbf{\textcolor{good}{0.800}} \\
child median & 0.783 & \textbf{\textcolor{good}{0.802}} \\
seed (executor re-score) & 0.746 & 0.753 \\
\bottomrule
\end{tabular}
\caption{\armB{} leads on all three central-tendency measures. The gap is modest ($\sim$0.02 in
mean) but consistent, and in the \emph{same direction} as the higher valid yield: constrained
edits stay close to a working retrieval chain, so the child distribution degrades less. Caveat:
one replicate per arm and no same-config variance floor on this task yet, so the $\sim$0.02
mean gap should be read as \emph{suggestive}, not a measured win; the validity gap, by contrast,
is categorical.}
\end{table}

\subsection{Held-out test coverage vs.\ the CARL feedback-condition baselines}
\label{sec:benchcmp}
The four inference-time feedback conditions (\emph{without-feedback}, \emph{self-reason},
\emph{metric-aware}, \emph{data-aware}) are an \textbf{executor} axis, orthogonal to this
report's \textbf{mutator} axis. Our run holds the executor fixed at \texttt{feedback\_mode=none}
(\emph{without-feedback}: single-pass, no retry) and varies only the mutation operator. To place
our winners on the same HoVer test set, Table~\ref{tab:benchcmp} reproduces the feedback-condition
sweep (mutator \texttt{Qwen3-235B-A22B-Thinking}) and adds our two best programs (mutator
\texttt{gemini-3.5-flash}). Metric is \textbf{discrete coverage} ($1$ iff every gold title
retrieved), $n=300$ test claims, base $=$ seed chain $\to$ best $=$ evolved winner.

\begin{table}[H]\centering\small
\begin{tabular}{lllcl}
\toprule
method & mutator & feedback & test disc.\ base$\to$best & source \\
\midrule
prompt & Qwen3-235B-Think & ---          & $43.3 \to 49.0$ & prior table \\
GEPA   & Qwen3-235B-Think & ---          & $43.3 \to 46.0$ & prior table \\
CARL   & Qwen3-235B-Think & none         & $48.7 \to 56.0$ & prior table \\
CARL   & Qwen3-235B-Think & self-reason  & $46.3 \to 57.0$ & prior table \\
CARL   & Qwen3-235B-Think & metric-aware & $46.3 \to 60.7$ & prior table \\
CARL   & Qwen3-235B-Think & data-aware   & $68.3 \to 68.7$ & prior table \\
\midrule
\armA{} (free rewrite)   & gemini-3.5-flash & none & $44.0 \to \textbf{63.7}$ & \textbf{this run} \\
\armB{} (CARL tool diff) & gemini-3.5-flash & none & $44.0 \to \textbf{\textcolor{good}{68.0}}$ & \textbf{this run} \\
\addlinespace
\armA{} (free rewrite)   & Qwen3-235B-Think & none & $44.0 \to$ \emph{pending} & running \\
\armB{} (CARL tool diff) & Qwen3-235B-Think & none & $44.0 \to$ \emph{pending} & running \\
\bottomrule
\end{tabular}
\caption{Held-out HoVer discrete coverage ($n=300$). \textbf{Top block:} the feedback-condition
sweep under the \texttt{Qwen3-235B-Thinking} mutator (user-supplied). \textbf{Bottom block:} this
run's two winners under the \texttt{gemini-3.5-flash} mutator, both at \texttt{feedback=none}. Our
seed re-scores to $44.0$, inside the prior table's base cluster ($43.3$--$48.7$) --- the two
evaluations measure the \emph{same} discrete metric on the same split. Two reads: (i) at the
\emph{same} without-feedback setting, \armB{} climbs $44.0\!\to\!68.0$ ($+24$) versus the Qwen
without-feedback row's $48.7\!\to\!56.0$ ($+7.3$) --- a \textbf{mutator-quality} gap, since the
mutator is the only difference; (ii) \armB's $68.0$ matches the strongest Qwen condition
(\emph{data-aware}, $68.7$), but that row's base is \emph{already} $68.3$ --- its score is an
inference-time effect (the ground-truth-hint retry lifts even the un-optimized seed), whereas
\armB{} reaches the same place by \emph{search alone}. The apples-to-apples control is a full
paired run (\textbf{both} arms under the same Qwen mutator, last two rows) --- it will show both
whether the Qwen mutator reproduces $\sim$56 or our $\sim$68, and whether the diff still beats
free-rewrite when the mutator itself is strong (early live fitness has the two arms near-parity). \emph{Caveats:} the source table's column semantics (Val vs Test) and the
identity of its ``CARL'' operator are our best-judgment reading; discrete-coverage definition
assumed identical across evaluations.}
\label{tab:benchcmp}
\end{table}

\subsection{Token cost --- a reasoning-instrumentation gap}
\label{sec:tokens}
The prior non-reasoning run cleanly showed the diff's output economy because reasoning tokens
were 0 on both arms. Here the split is only half-recorded, and honesty requires we say why: the
\texttt{LLM\_CALL} event emitted by \armA's \texttt{MutationAgent} (and the base agent) never
forwarded the \texttt{tokens\_reasoning} field, while \armB's \texttt{DiffMutationAgent} did. So
the ``0'' for \armA{} below is \emph{not recorded}, not measured-zero --- both arms sent
\texttt{reasoning.effort:high} and gemini almost certainly reasoned on both. Only
\texttt{tokens\_out} (total generated, which \emph{includes} reasoning) is logged correctly on
both arms; it is the only cross-arm comparison this run supports.
\begin{table}[H]\centering
\begin{tabular}{lccc}
\toprule
mutation side (\texttt{gemini-3.5-flash}), full run & \armA{} & \armB{} & \armB{} vs \armA{} \\
\midrule
tokens in & 1{,}715{,}634 & 1{,}934{,}479 & $+12.8\%$ \\
\textbf{tokens out (total)} & \textbf{1{,}678{,}069} & \textbf{1{,}882{,}208} & \textbf{$+12.2\%$} \\
\quad of which reasoning (recorded) & \textcolor{bad}{not recorded$^\dagger$} & 1{,}639{,}000 (87\%) & --- \\
\quad content (out $-$ recorded reasoning) & \textcolor{bad}{n/a$^\dagger$} & 243{,}208 & --- \\
LLM calls (incl.\ retries) & 252 & 259 & $+3\%$ \\
\bottomrule
\end{tabular}
\caption{$^\dagger$\armA's reasoning tokens were dropped by the \texttt{MutationAgent}
\texttt{LLM\_CALL} emitter (a pre-existing bug, since fixed on both emit paths), so \armA's
reasoning/content split is \emph{unrecoverable for this run} --- \armA{} is unmeasured-reasoning,
not zero-reasoning. The only clean number is \textbf{total output}: \armB's is
\textbf{+12\% higher}. By construction the diff emits slot deltas rather than a whole
re-serialised tool chain, so its \emph{content} payload is genuinely small (243k tokens, 13\% of
its output) --- but we cannot state the rewrite-vs-diff content ratio this run, because \armA's
content is entangled with its unrecorded reasoning. On this run the diff's demonstrated advantage
is validity, fitness, and legibility; the total-output comparison is the one token fact that
holds, and on it the diff is dearer.}
\end{table}
\begin{figure}[H]\centering
\includegraphics[width=\textwidth]{ab_tokens.png}
\caption{Left: mutation output tokens per call (dots) with a 10-call rolling mean --- the two
arms overlap because \armB's small content payload plus its reasoning channel roughly matches
\armA's output (content plus its \emph{unrecorded} reasoning; see Sec.~\ref{sec:tokens}). Right:
cumulative tokens --- \armB's input (dashed) and output (solid) both sit slightly above \armA's;
the curves do not separate the way they did on the non-reasoning run.}
\end{figure}

\subsection{Cost/quality frontier --- cost-to-winner, not total spend}
\label{sec:pareto}
The table above counts \emph{tokens}; converting them to dollars at the cheapest OpenRouter
\texttt{gemini-3.5-flash} price (\$1.50/M in, \$9.00/M out, \emph{output inclusive of reasoning};
the Qwen3-8B executor is a separate, shared cost and excluded here) reframes the comparison ---
and the answer flips depending on \emph{which} dollar you count. Cost is attributed by
\texttt{created\_at}: each child, in wall-clock creation order, consumes the next accepted
mutation call, so ``cost to a program'' is the cumulative mutation spend up to and including that
program's creation. (We deliberately do \emph{not} use the persisted \texttt{atomic\_counter}
field --- it is non-monotone against creation time here and mis-orders the trajectory; \texttt{created\_at}
and the evolution-loop \texttt{iteration} clock agree to within a rank or two throughout.)

\begin{table}[H]\centering
\begin{tabular}{lccc}
\toprule
mutation-model \$ (OpenRouter cheapest) & \armA{} & \armB{} & \armB{} vs \armA{} \\
\midrule
total run spend & \$17.68 & \$19.84 & $+12\%$ \\
\quad of which reasoning channel & \textcolor{bad}{not itemised$^\dagger$} & \textcolor{bad}{\$14.75 (74\%)} & --- \\
\midrule
\textbf{cost to the winning program} & \$15.93 & \textbf{\textcolor{good}{\$6.46}} & \textbf{$-59\%$} \\
\quad winner reached at (chrono.\ rank) & 223/250 (89\%) & \textbf{\textcolor{good}{80/254 (31\%)}} & far earlier \\
\quad winner fitness (train soft) & 0.831 & \textbf{\textcolor{good}{0.858}} & higher \\
\midrule
cost per \emph{valid} program & \$0.119 & \textbf{\textcolor{good}{\$0.083}} & $-30\%$ \\
budget wasted on invalid children & \textcolor{bad}{\$6.8 (38\%)} & \textbf{\textcolor{good}{\$0.2 (1\%)}} & --- \\
\bottomrule
\end{tabular}
\caption{Two dollars, two answers. On \textbf{total} spend the diff is \emph{dearer} ($+12\%$).
$^\dagger$\armB's recorded reasoning channel is \$14.75 (74\% of its \$19.84); \armA's reasoning
was never itemised (Sec.~\ref{sec:tokens}) so we cannot net it out --- \armA{} \emph{was} billed
for its full generated output, which includes its own unrecorded reasoning. So the total-spend
gap is real but its reasoning attribution is one-sided. On the dollar that
actually matters for a search that plateaus early --- spend until the best program exists ---
\armB{} buys a \emph{higher}-fitness winner (0.858 vs 0.831) for \textbf{\$6.46 vs \$15.93}, a
$59\%$ saving. The reason is timing: \armB's winner appears at chronological rank 80 of 254
(31\% of the run) and is \emph{never beaten}, so \armB{} could have stopped at \$6.46 and kept its
best; \armA's winner does not arrive until rank 223 of 250 (89\%), so almost the whole budget is
load-bearing. Per \emph{usable} program the diff is cheaper outright (\$0.083 vs \$0.119),
because \armA{} burns \$6.8 (38\% of its spend) on the 93 malformed/invalid children that never
enter selection while \armB{} burns \$0.2 (1\%) on 3.}
\end{table}
\begin{figure}[H]\centering
\includegraphics[width=\textwidth]{ab_pareto.png}
\caption{\textbf{Left:} cost-to-winner vs.\ best train fitness. The step line is each arm's
cost-efficiency path (best-so-far fitness as cumulative mutation dollars accrue); \armB{} sits
up-and-to-the-left --- higher fitness reached for fewer dollars --- and owns the Pareto frontier.
\textbf{Right:} cost-to-program vs.\ held-out test discrete coverage, every evaluated lineage
program placed at its true cost-up-to-creation. \armB's connected climb reaches $68\%$ coverage
by \$6.46, with the decisive \texttt{gen1}$\to$\texttt{gen2} rewrite already at $64.3\%$ for a
cumulative \textbf{\$0.69}; \armA's free-rewrite winner reaches only $63.7\%$ and not until \$15.9.
x-axis is mutation-model tokens only (OpenRouter cheapest, reasoning included); executor excluded.}
\end{figure}
\begin{figure}[H]\centering
\includegraphics[width=\textwidth]{ab_waste.png}
\caption{\textbf{Left:} where each arm's mutation budget went, split (by the arm's valid/invalid
child rate) into productive spend on valid children (green) and spend wasted on malformed/invalid
children (red). \armA{} throws away \$6.8 of \$17.68 ($38\%$) on children that never enter
selection; \armB{} wastes \$0.2 of \$19.84 ($1\%$). \textbf{Right:} effective cost per \emph{valid}
program --- \$0.083 for the diff vs \$0.119 for free rewrite, a $30\%$ reduction despite \armB's
higher total spend, because almost all of \armB's dollars land on programs that count.}
\end{figure}

\subsection{Structural exploration --- both preserve topology, \armB{} tighter}
\label{sec:struct}
On this task the operators diverge less than on the summarizer run: both mostly preserve the
seed's 7-step structure. \armA's valid children are $\{7\!:\!115,\ 6\!:\!14,\ 5\!:\!20\}$ (77\%
at 7 steps), spreading modestly toward shorter chains; \armB's are
$\{7\!:\!200,\ 6\!:\!20,\ 5\!:\!18,\ 3\!:\!1\}$ (84\% at 7 steps), with a single deliberate
$7\!\to\!3$ prune. There is no wholesale chain-collapse here (the retrieval task rewards keeping
tool steps, so even the free rewriter keeps them), but the diff stays marginally closer to the
parent topology --- consistent with edits that touch slots rather than re-imagine the chain.
\begin{figure}[H]\centering
\includegraphics[width=0.72\textwidth]{ab_nsteps.png}
\caption{Chain-length distribution of valid children. Both arms concentrate at 7 steps; \armA{}
spreads a little more toward 5--6, \armB{} stays tighter at 7 with one 3-step prune.}
\end{figure}

\section{The winning chains --- convergent architecture, divergent reliability}
\label{sec:chains}
The two arms produced \emph{step-for-step the same} best retrieval architecture. Both winners
abandon the seed's ``retrieve $\to$ summarise $\to$ query'' rhythm and converge on a
\textbf{deepen-and-accumulate} chain: four \texttt{retrieve\_deep} ($k{=}10$) tool steps
alternating with three query-generation LLM steps, and each query-generator fed the first-hop
retrieval plus its immediate predecessor (\texttt{deps} $[1,3]$ then $[1,3,5]$). The summarise
steps --- an information bottleneck that discarded raw passages before the next hop --- are gone
in both.

\begin{table}[H]\centering\small
\begin{tabular}{clll}
\toprule
step & type & \armA{} best (0.831) & \armB{} best (0.858) \\
\midrule
1 & \textcolor{accent}{tool} & \texttt{retrieve\_deep(\$outer)} & \texttt{retrieve\_deep(\$outer)} \\
2 & \textcolor{bad!70!black}{llm} & 2nd-hop query \hfill $[1]$ & 2nd-hop query \hfill $[1]$ \\
3 & \textcolor{accent}{tool} & \texttt{retrieve\_deep(\$hist[-1])} \hfill $[2]$ & \texttt{retrieve\_deep(\$hist[1])} \hfill $[2]$ \\
4 & \textcolor{bad!70!black}{llm} & 3rd-hop query \hfill $[1,3]$ & 3rd-hop query \hfill $[1,3]$ \\
5 & \textcolor{accent}{tool} & \texttt{retrieve\_deep(\$hist[-1])} \hfill $[4]$ & \texttt{retrieve\_deep(\$hist[3])} \hfill $[4]$ \\
6 & \textcolor{bad!70!black}{llm} & 4th-hop query \hfill $[1,3,5]$ & 4th-hop query \hfill $[1,3,5]$ \\
7 & \textcolor{accent}{tool} & \texttt{retrieve\_deep(\$hist[-1])} \hfill $[6]$ & \texttt{retrieve\_deep(\$hist[5])} \hfill $[6]$ \\
\bottomrule
\end{tabular}
\caption{The two winners are the same chain. Identical step types, identical topology, identical
dependency sets. The \emph{only} difference is how a tool step names its query source: \armA{}
uses the relative \texttt{\$history[-1]} (last output), \armB{} the absolute positional
\texttt{\$history[1]/[3]/[5]} --- functionally equivalent on a linear chain. This is convergent
evolution: both operators find the ``deepen every hop, drop the summariser'' architecture. The
diff's edge is not a \emph{different} destination but the \emph{reliability of the path} to it ---
\armB{} reached this design having malformed \textbf{zero} children en route, \armA{} while
malforming 38\%.}
\end{table}

\subsection{How \armB{} evolved it: the base\,$\to$\,best lineage}
\label{sec:lineage}
The best \armB{} child sits at the end of a clean four-mutation lineage from the seed. Every
child is a structured slot-diff of its parent, so the whole path is legible as a sequence of
typed edits --- exactly what the free-rewrite arm cannot offer, since each of its children is an
independent whole-document re-emission with no diff to inspect.

\begin{figure}[p]\centering
\includegraphics[width=\textwidth,height=0.92\textheight,keepaspectratio]{lineage_graph.png}
\caption{\armB's winning lineage as a \textbf{dependency graph}, one column per program
(base\,$\equiv$\,gen1\,$\to$\,gen4). Each column is the seven-step pipeline drawn top-to-bottom;
every node is a step, colored by type (\textcolor{accent}{blue}~=~\texttt{retrieve}/%
\texttt{retrieve\_deep} tool, \textcolor{bad!70!black}{orange}~=~LLM) and labelled with its role.
\textbf{Arrows are the step \texttt{dependencies}} --- the wiring that defines the DAG. An edge is
grey when it is unchanged from the parent, \textcolor{good}{solid green} when the mutation
\emph{added} it, and \textcolor{bad}{dashed red} when it \emph{removed} it; a node is outlined in
\textcolor{bad}{red} when any of its fields changed versus the previous generation. The italic line
under each new generation is the CARL structured-diff operator's own stated change --- its words,
not a reconstruction. Read left to right, the graph \emph{is} the story: \textbf{gen2} fans green
skip-edges into the two query steps (\#4,\,\#6) so each hop's query sees the whole accumulated
history --- the decisive edit ($\Delta$fit $+0.094$), a coordinated seven-step rewrite where every
tool step also becomes \texttt{retrieve\_deep} and the summarise steps flip to query-generators;
\textbf{gen3} strips those fan-ins back (dashed red) to a clean linear chain while rewriting the
reasoning prompts ($+0.009$); \textbf{gen4} re-adds a \emph{selective} subset of skip-connections
in green for the final $+0.002$. \textbf{base\,$\equiv$\,gen1} carries no red --- the first mutation
was a representable no-op. Because each child is a diff against a known parent, the model touches
only the wiring it means to --- never re-emitting (and re-risking) the untouched steps, which the
free-rewrite arm cannot avoid.}
\label{fig:lineage}
\end{figure}

\begin{table}[H]\centering\small
\begin{tabular}{llccc}
\toprule
program & structural change vs.\ parent & fitness (train soft) & test soft & test disc.\ (\%) \\
\midrule
base (seed)     & ---                              & ---   & 0.752 & 44.0 \\
gen 1           & no-op (structurally identical)   & 0.753 & 0.763 & 47.7 \\
gen 2           & rewire every hop; \texttt{+}skip-edges & 0.847 & 0.848 & 64.3 \\
gen 3           & prune fan-ins to linear chain    & 0.856 & 0.859 & 66.7 \\
gen 4 (best)    & re-add selective skip-edges      & \textbf{0.858} & \textbf{0.862} & \textbf{\textcolor{good}{68.0}} \\
\midrule
\armA{} best (free rewrite) & --- (no inspectable diff) & 0.831 & 0.839 & 63.7 \\
\bottomrule
\end{tabular}
\caption{Held-out test payoff of \emph{every} program on the lineage ($n=300$ test claims; fitness
is the train soft-coverage objective, 300 train claims). Held-out discrete coverage climbs
monotonically down the four-mutation path ($44.0\to47.7\to64.3\to66.7\to\textbf{68.0}$) --- the
structural edits shown in Fig.~\ref{fig:lineage} each carry a real test gain, not just a train one,
and soft coverage tracks it in lockstep. base and gen1 are the \emph{same} structure yet score
$44.0$ vs $47.7$: the $\sim$3.7-pt gap is the non-deterministic \texttt{Qwen3-8B} executor's noise
floor, the band against which the per-generation deltas should be read. The evolved \armB{} winner
($68.0$) clears the \armA{} free-rewrite winner ($63.7$) on the same held-out split by $4.3$ pts.}
\label{tab:lineagetest}
\end{table}

\noindent Two structural observations. First, the operator can emit a \emph{no-op} diff
(\emph{base\,$\to$\,gen1}): keeping every slot unchanged is a representable, valid move --- wasted
compute, but never a malformed child. Second, the load-bearing mutation (\emph{gen1\,$\to$\,gen2})
is a \emph{coordinated multi-field, multi-step rewrite} --- tool swaps, step-type flips, and
dependency rewiring together --- yet it is expressed as one structured diff that the compiled
schema guarantees is well-formed. The very edit that is most likely to corrupt a free-form
re-emission (touch every step at once) is exactly where the diff's constraints pay off.

\section{Engineering deliverable: per-candidate executor token accounting}
\label{sec:tokentrack}
The chain executor (\texttt{Qwen3-8B}) already captured per-call \texttt{prompt\_tokens}/%
\texttt{completion\_tokens} in its \texttt{CallLog}, but the chain runners call
\texttt{client.copy()} for every concurrent LLM call and the base \texttt{copy()} hands out a
\emph{fresh} log list for thread-safety --- so usage landed in throwaway copies and the parent's
logs stayed empty, and \texttt{validate()} discarded it. This run therefore has \emph{no}
per-candidate executor tokens; the only executor figure available is a shared-proxy aggregate:
\textbf{$\approx$132M tokens combined across both arms} (proxy Prometheus counter at
\texttt{10.232.89.98:4000}, no per-arm split and shared with any co-tenant traffic).

That gap is now closed. \texttt{problems/chains/usage.py} adds a \texttt{LogAggregatingLLMClient}
whose \texttt{copy()} shares one log list across the fan-out, plus \texttt{usage\_totals()}
(prompt/completion/total tokens, \texttt{n\_llm\_calls}, \texttt{llm\_cost}) and a
\texttt{ZERO\_USAGE} constant. It is wired into \texttt{full7/validate.py} (success and
early-reject paths) and refactors the summarizer validator onto the same helper; four unit tests
cover the copy-sharing invariant. \textbf{Future} HoVer runs will carry per-candidate executor
tokens in program metrics, visible via \texttt{gigaevo top} and the A/B analyzer --- enabling
the per-arm executor-cost comparison this run could only approximate.

\section{Verdict}
The structured diff generalises cleanly to tool-bearing retrieval chains and the validity
guarantee holds --- decisively --- even against a strong reasoning mutator. Over completed
evaluations, free-form rewriting yielded \textbf{61.6\% valid (91/242 malformed JSON, 38\%)}
while the tool diff yielded \textbf{98.8\% valid (0 malformed)}; that gemini's malformed rate
matches the weak model's says the failure is structural to full-document re-emission, not model
weakness. \armB{} also leads fitness (mean 0.800 vs 0.778, best 0.858 vs 0.831) and keeps chains
tighter to the seed topology. On cost the picture is two-sided: a reasoning-instrumentation gap
(\armA's \texttt{MutationAgent} \texttt{LLM\_CALL} emitter dropped the reasoning field, so only
\armB's 87\% reasoning channel was recorded --- \armA's is unmeasured, not zero; bug now fixed)
leaves \emph{total} billed output as the only clean cross-arm read, and on it \armB's is 12\%
higher --- so on total spend the diff is dearer. But
on the dollars that matter for the search, the diff wins outright: priced at OpenRouter's cheapest
\texttt{gemini-3.5-flash} rate it reaches a \emph{higher} winner (0.858 vs 0.831) for \textbf{\$6.46
vs \$15.93} cost-to-winner ($-59\%$, its winner arriving at 31\% of the run vs 89\% and never
beaten) and costs \textbf{\$0.083 vs \$0.119} per valid program ($-30\%$), because free rewrite
wastes 38\% of its budget on malformed children the diff never emits. Caveats:
single replicate per arm, no same-config variance floor (so the fitness gap is suggestive, the
validity gap categorical), \texttt{memory=none} (citation-integrity instrumented but trivially
0/0), \texttt{num\_parents=1} (no crossover), and 14 in-flight children per arm excluded from
completed figures. Next: a variance floor on HoVer full7; a memory-on / two-parent run to
exercise crossover tool-diffs and the evidence surface; and a re-run once
\texttt{problems/chains/usage.py} is live to settle the per-arm executor-token question and,
with the reasoning field now emitted on both arms, to record a clean per-arm reasoning split.
\end{document}
